\chapter{Planning and Conducting Yoga Lessons}
\label{ch:lessons}

\emph{How does an instructor plan a lesson, adapt it to the group in front of them, run the class, keep it safe, and carry themselves as a teacher?}
The teaching skills unit asks for the essentials of a lesson plan, the principles of teaching different groups, class management, precautions and contraindications, the features of an ideal instructor, and model lesson plans. The preparation for a class and the factors that influence teaching are in \cref{ch:teaching-class}; this chapter covers the rest.

\begin{unverified}
This chapter is written from notes blocks whose fact check is still pending. Points follow the Board syllabus and guideline, the Common Yoga Protocol and the official guidebook's chapter on teaching practice, as marked in the notes; the remaining points are the author's workshop teaching. A certified teacher must review it before print.
\end{unverified}

\section{Essentials of a good lesson plan}

\subsection*{Why teach yoga, and why plan}
Teaching yoga spreads health and happiness, helps the teacher to grow, and opens many kinds of work. A lesson plan makes the class safe, balanced and on time: it states the aim, the practices, their order, the time for each and the cautions. It is a short written page, not an essay.

\subsection*{Before planning}
Find out what the learners need and how much experience they have. The guidebook asks for a student-centred approach, with sessions fitted to the level of the students and objectives, activities and assessment that line up.

\subsection*{The parts of a plan}
A plan names the group (age, level, health), the aim, the practices in order, the time for each, the instructions, and the closing. The Board names four areas to plan for: the cleansing practice, asana, pranayama, and the practices leading to dhyana. The plan is planned but flexible: the guidebook asks the teacher to change the method to suit the learners.

\subsection*{The order of a class}
The Board's own order, in the common protocol, is: prayer, loosening practices, asanas (standing, sitting, prone, then lying on the back), kapalabhati, pranayama, dhyana, sankalpa, and the closing peace chant, with relaxation as needed. The workshop summarises this as warm-up, main practice and cool-down; the guidebook and protocol do not use those words.

\subsection*{Eight steps to teach one asana}
The workshop teaches one asana in eight steps. Schools use their own step lists; keep this as the exam form and fit it to the Board protocol.
\begin{enumerate}
\item Give the name and its meaning, with the reason for the name and the category of the posture.
\item Show the starting posture and the resting posture.
\item State the limits (who should not do it) and the benefits.
\item Demonstrate; then the class practises.
\item Practice in pairs: one group practises while the other watches and corrects, then they swap.
\item The demonstrator repeats, stressing the key points (knees, eyes, breath).
\item The whole group practises together.
\item Questions and answers, then the closing prayer.
\end{enumerate}

The starting and resting postures by group are: standing, tadasana and a loose standing rest; sitting, dandasana and a loose sitting rest; prone, the lying start and makarasana; supine, the lying start and shavasana. The class traditionally closes with the salutation to Patanjali, a prayer for a healthy mind, speech and body.

\section{Teaching different groups}

The Board syllabus names six groups: beginners, children, youth, women, the geriatric population and special attention groups. Say how the practice is adapted for each, not what it cures.
\begin{itemize}
\item \textbf{Beginners} need a general background in yoga, motivation and the habit of practice: short, simple practices in plain words. An experienced group wants more practices and finer detail, longer holds and subtler work. The texts speak of three levels of aspirant, beginner, middle and advanced, in the same spirit.
\item \textbf{Children} of school age, about 6 to 16, need short, lively and interactive practices fitted to their age and temperament, with their attention span in mind.
\item \textbf{Youth} have more energy and range: practices that build strength, focus and balance (the author's wording; the guidebook gives no separate guidance).
\item \textbf{Women}: take note of the menstrual cycle, pregnancy and the body's needs, and avoid the practices cautioned for them. In pregnancy and during menstruation, a yoga expert is asked first.
\item \textbf{Older adults} need slower, supported practice, seated or against a wall, with gentle movement and shorter holds (the author's wording).
\item \textbf{Special attention groups}, people with a health condition or a disability, are taught in a group of similar people where possible, and only within medical advice.
\end{itemize}

\ycbfig{ycb1-fig-teaching-group}

\section{Class management}

\subsection*{Meaning and need}
Class management is the organising of the people, the space, the time and the attention so that the class runs smoothly and safely. It keeps the practice safe, uses the time well, and keeps the students calm and interested. Keep the class small enough to see everyone.

\subsection*{Giving instructions}
Get everyone's attention first. Make the instruction short, clear and loud enough, one thing at a time, and describe a practice before it starts. If it is not understood, say it in other words. Explain long points while the group sits near and rests between two practices, never in the middle of a posture; a question-and-answer session closes a unit.

\subsection*{Demonstration}
Show the practice first, then say the key points. The students follow, and the teacher corrects with care. The guidebook teaches spoken instruction and demonstration as the two main methods and puts demonstration before the students' practice.

\subsection*{Questions and conduct}
Invite questions and never mock one. Keep the talk to the point; if you do not know an answer, say you will find out. Begin the session in silence with a meditative practice; teach the limits of a practice before introducing it; and see that everyone wears loose, comfortable clothes.

\subsection*{Time and difficulty}
Plan the minutes, start and end on time, and put the important practices first. When a student finds a posture hard, modify the posture or its alignment quickly and build trust; encourage patience and regular practice within the student's capacity. Never force the body, and never ignore the student.

\ycbfig{ycb1-fig-teaching-lecture}

\section{Precautions and contraindications}

\subsection*{Principles of the practical lesson}
Yoga is not vigorous exercise. Practise slowly, with the breath, within your own capacity; hold a posture with little effort; finish with relaxation; rest in shavasana if tired.

\subsection*{Five working rules}
\begin{enumerate}
\item Assess each person's health and physical limits before the class.
\item Modify postures to fit the person.
\item Watch the students for discomfort or strain.
\item Do not let anyone overexert.
\item Tell the class the contraindications of each practice before it is taught.
\end{enumerate}

\subsection*{Safety}
Correct by word and demonstration first. Touch only with the student's consent and with care. Stop a student who feels pain, dizziness or breathlessness. Do not let anyone practise during an illness or fever, or after surgery, a fracture or a sprain, without a physician's advice. Each practice has its own cautions (high blood pressure, heart disease, hernia, back problems, pregnancy, menstruation and others); they are given with each practice in the practical chapters, and the teacher keeps them in mind.

\subsection*{Scope and referral}
The instructor teaches the common protocol for prevention and well-being. Do not diagnose and do not promise a cure. Refer a person to a doctor or another teacher when you cannot give them enough time, when the need is outside your skill, or when they have gone as far as you can help.

\begin{caution}
A good lesson is safe first: ask about health, show, correct kindly, and stop at pain.
\end{caution}

\ycbfig{ycb1-fig-teaching-demo}

\section{The ideal yoga instructor}

\subsection*{Qualities the Board names}
The Board's guideline lists the attributes for this level: self-discipline, confidence, maturity, patience, compassion, active listening, time management, empathy and proficiency in language, with good communication and the ability to read the trainees' body language. The workshop's five words, knowledgeable, communicative, empathetic, professional and adaptable, are the author's summary, not the Board's list.

\subsection*{Skills}
Prepare well, demonstrate clearly, explain simply, and adapt the practice for each student. The guidebook expects the teacher to know more than the students and to be able to modify a practice.

\subsection*{Communication}
Say what you mean in few words. Ask questions to be sure you are understood, and listen without hurry. Give more appreciation than complaint, turn a criticism into a polite request, and do not blame.

\subsection*{Conduct}
Be punctual and calm; welcome every student; respect their privacy and their limits; and keep learning and practising yourself. Be ready to meet misunderstanding with courage and a steady mind.

\ycbfig{ycb1-fig-teaching-teacher}

\section{Model lesson plans}

The model order is the protocol's: prayer, loosening practices, asanas (standing, sitting, prone, supine), kapalabhati, pranayama, meditation, sankalpa and the closing peace chant. Two adaptations show the method:
\begin{itemize}
\item \textbf{A mixed beginners' group}: a short prayer; loosening practices from the neck down; a few easy asanas; relaxation; one or two gentle breathing practices; a short meditation; the closing prayer.
\item \textbf{A special group}: fewer and slower practices, in seated or lying positions, with more rest, and no practice cautioned for the condition; all within medical advice.
\end{itemize}
To adapt, shorten or lengthen each part, keep the order, and drop what is not suitable. A good model plan follows the Board's order, matches the aims to the learners, gives each part its time, includes a range of asana, pranayama and meditation, and stays flexible. Sample plans with minutes are still to be written, once the protocol's timings are checked; the guidebook and protocol give no minutes for a class.

\ycbfig{ycb1-fig-teaching-lesson}

\begin{recall}
\begin{itemize}
\item Plan: group, aim, practices in order, minutes, instructions, cautions, closing.
\item Protocol order: prayer, loosening, asanas, kapalabhati, pranayama, dhyana, sankalpa, peace chant.
\item Groups: beginners, children, youth, women, geriatric, special attention.
\item Instructions: attention first, short and clear, one at a time; explanations at rest.
\item Safety: ask, show, correct kindly, stop at pain; refer beyond your scope.
\end{itemize}
\end{recall}

\subsection*{Worked examples}
\begin{itemize}
\item \textbf{Model answer: a plan for older adults in a community hall.} Prayer; seated neck, shoulder and ankle loosening; tadasana with a chair for support; seated vakrasana; lying pavanamuktasana with one leg; shavasana; abdominal breathing and nadi shodhana without holding; a short meditation; sankalpa and the peace chant. Leave out forceful breathing and deep back bends.
\item \textbf{A student keeps asking questions in the middle of a posture.} Ask them to hold the question for the rest that follows, and answer it there for everyone.
\item \textbf{A student says their back hurts in bhujangasana.} Stop them, have them rest in makarasana, offer sarala bhujangasana or another posture, and suggest they consult a doctor if the pain continues.
\end{itemize}

\subsection*{Check yourself}
\begin{enumerate}
\item What should a teacher find out before planning?
\item Give the class order of the common protocol.
\item Name the starting and resting postures for standing asanas.
\item What happens in the step ``practice in pairs''?
\item Should the day's lesson plan be rigid?
\item How is practice adapted for older adults?
\item When are questions best taken in a group class?
\item When should the limits of a practice be taught?
\item When should a student not practise without a doctor's advice?
\item What should the instructor do when a need is beyond their scope?
\item Name three personal qualities the Board lists for this level.
\end{enumerate}

\subsection*{Answers}
\begin{enumerate}
\item The learners' needs and level of experience.
\item Prayer, loosening practices, asanas, kapalabhati, pranayama, dhyana, sankalpa, peace chant.
\item Tadasana and a loose standing rest.
\item One group practises while the other watches and corrects, then they swap.
\item No; it is planned but flexible.
\item Slower, supported (seated or against a wall), gentle movement and shorter holds.
\item During the rest between two practices.
\item Before the practice is introduced.
\item During an illness or fever, or after surgery, a fracture or a sprain.
\item Refer the person to a doctor or another teacher.
\item Any three of: self-discipline, confidence, maturity, patience, compassion, active listening, time management, empathy, language proficiency.
\end{enumerate}

\paragraph{Sources.} The Board syllabus for Level 1 (teaching skills unit) and the Board Guideline version 0.3 (attributes); the Common Yoga Protocol, Ministry of AYUSH (class order); the official guidebook for Level I (2016, older scheme), chapter on teaching practice; the author's workshop teaching for the eight steps and the notes marked as the author's wording.
